\begin{statementbox}

	\textbf{\textcolor{CStatement}{条件}}
	\begin{description}[style=nextline, leftmargin=2.8em, labelsep=0.5em, itemsep=0.35em]
		\item[\textbf{\textcolor{CStatement}{条件 1（全集给定）：}}] 设全集为 $U$，集合 $A$ 满足 $A \subseteq U$。
	\end{description}

	\textbf{\textcolor{CStatement}{结论}}
	\begin{description}[style=nextline, leftmargin=2.8em, labelsep=0.5em, itemsep=0.45em]
		\item[\textbf{\textcolor{CStatement}{结论一（双重补集）：}}]
			\textbf{\textcolor{CStatement}{直观描述：}}
			一个集合补两次，结果回到自身。
			\[
				\complement_U(\complement_U A) = A
			\]

		\item[\textbf{\textcolor{CStatement}{结论二（边界补集）：}}]
			\textbf{\textcolor{CStatement}{直观描述：}}
			全集的补是空集，空集的补是全集。
			\[
				\complement_U U = \varnothing,
				\qquad
				\complement_U \varnothing = U
			\]

		\item[\textbf{\textcolor{CStatement}{常见变形（差集形式）：}}]
			\textbf{\textcolor{CStatement}{直观描述：}}
			引入差集符号，双补还原可写为更熟悉的形式。
			\[
				\complement_U A = U \setminus A
			\]
			于是双重补集还原律等价于
			\[
				U \setminus (U \setminus A) = A
			\]
	\end{description}

	\textbf{\textcolor{CStatement}{使用提醒（全集依赖）：}}
	补集符号 $\complement_U A$ 必须明确全集 $U$；全集改变，补集结果随之改变。

\end{statementbox}
